No treinamento supervisionado de modelos de Inteligência Artificial, em especial de
\glspl{rna}, o aprendizado ocorre pelo ajuste iterativo dos parâmetros internos da rede
(pesos e vieses) a partir de exemplos rotulados \cite{jesus2023inteligencia}. Para que
esse ajuste seja possível, é necessário quantificar matematicamente o quanto a saída
prevista pelo modelo diverge do resultado esperado. Essa medição é realizada pela
\textbf{função de perda} (\textit{loss function}), também denominada função de custo ou
função objetivo.

Formalmente, dado um conjunto de entradas associado aos seus respectivos gabaritos $y$ e às
saídas estimadas $\hat{y}$ produzidas pela rede parametrizada por $\theta$, a função de
perda $\mathcal{L}(y, \hat{y}; \theta)$ atribui um valor escalar não negativo ao erro
cometido na previsão. Quanto maior a discrepância entre $\hat{y}$ e $y$, maior o valor
retornado pela função de perda. O treinamento consiste, portanto, em um problema de
otimização no qual se busca minimizar o valor médio da função de perda sobre os dados de
treinamento por meio de algoritmos baseados em gradiente, como a retropropagação do erro
(\textit{backpropagation}) associada a otimizadores como o Gradiente Descendente
Estocástico ou o Adam \cite{goodfellow2016deep}.

\subsubsubsection{Impacto do Desbalanceamento de Classes na Escolha da Perda}

A escolha da função de perda condiciona diretamente o comportamento do modelo durante a
convergência. Em problemas de segmentação semântica de imagens, como a delimitação de uma
\gls{roi} pela arquitetura U-Net, a rede classifica cada
pixel da imagem como pertencente ao objeto de interesse ou ao fundo. Uma função clássica
para essa tarefa é a \textbf{Entropia Cruzada Binária} (\textit{Binary Cross-Entropy} --
BCE), definida sobre $N$ pixels pela Equação~\eqref{eq:bce_loss}:

\begin{equation}
    \label{eq:bce_loss}
    \mathcal{L}_{\text{BCE}} = -\frac{1}{N} \sum_{i=1}^{N} \Big[ g_i \log(p_i) + (1 - g_i)\log(1 - p_i) \Big]
\end{equation}

\noindent em que $g_i \in \{0, 1\}$ representa o rótulo verdadeiro do $i$-ésimo pixel e
$p_i \in (0, 1)$ corresponde à probabilidade prevista pelo modelo para a classe positiva.

Embora a entropia cruzada avalie cada pixel de forma independente e favoreça boa
calibração probabilística, ela atribui o mesmo peso a todos os pixels da imagem. Quando a
região alvo ocupa apenas uma pequena fração do quadro (cenário típico da janela de dígitos
em imagens de medidores), ocorre um severo desbalanceamento entre a classe positiva
(objeto) e a classe negativa (fundo). Nessa condição, a rede pode atingir um valor baixo
de $\mathcal{L}_{\text{BCE}}$ simplesmente classificando quase todos os pixels como fundo,
o que compromete a detecção da estrutura de interesse \cite{jadon2020survey}.

\subsubsubsection{\textit{Dice Loss}}

Para contornar o desbalanceamento espacial em tarefas de segmentação sem exigir
reponderação manual das classes, \citeonline{milletari2016vnet} propuseram a
\textbf{\textit{Dice Loss}}, uma função de perda derivada do coeficiente de similaridade
de Sørensen--Dice, originalmente formulado para medir a sobreposição entre dois conjuntos
discretos. Como a interseção binária de conjuntos não é diferenciável (requisito
indispensável para o cálculo dos gradientes na retropropagação), utiliza-se na otimização
de redes neurais a sua variante contínua, denominada \textbf{perda \textit{Dice} suave}
(\textit{Soft Dice Loss}), expressa pela Equação~\eqref{eq:dice_loss}:

\begin{equation}
    \label{eq:dice_loss}
    \mathcal{L}_{\text{Dice}} = 1 - \frac{2 \sum_{i=1}^{N} p_i g_i + \epsilon}{\sum_{i=1}^{N} p_i + \sum_{i=1}^{N} g_i + \epsilon}
\end{equation}

\noindent em que $p_i \in [0, 1]$ é a probabilidade prevista para o pixel $i$,
$g_i \in \{0, 1\}$ é o valor correspondente na máscara de referência e $\epsilon > 0$ é
uma constante pequena de suavização adicionada ao numerador e ao denominador para evitar
divisão por zero e estabilizar o cálculo numérico quando a máscara alvo é vazia.

A formulação da \textit{Dice Loss} apresenta três propriedades relevantes para a
segmentação de regiões pequenas:

\begin{itemize}
    \item \textbf{Limitação de escala}: o termo fracionário varia entre $0$ (ausência
    total de sobreposição) e $1$ (sobreposição perfeita), fazendo com que
    $\mathcal{L}_{\text{Dice}} \in [0, 1]$ seja minimizada à medida que a máscara
    prevista converge para a máscara real.
    \item \textbf{Invariância ao tamanho do fundo}: diferentemente da entropia cruzada, o
    denominador da Equação~\eqref{eq:dice_loss} normaliza a interseção apenas pela soma
    das áreas previstas e reais da classe positiva, de modo que os pixels verdadeiros
    negativos do fundo não diluem o gradiente gerado pela região alvo.
    \item \textbf{Combinação com a entropia cruzada}: conforme analisado por
    \citeonline{jadon2020survey}, em estágios iniciais do treino, quando a sobreposição
    $p_i g_i$ ainda é muito baixa, os gradientes da \textit{Dice Loss} podem apresentar
    maior variância. Por essa razão, é prática recorrente combinar ambas as formulações em
    uma perda híbrida ($\mathcal{L}_{\text{total}} = \mathcal{L}_{\text{BCE}} +
    \mathcal{L}_{\text{Dice}}$), unindo a estabilidade de convergência pixel a pixel da
    entropia cruzada à sensibilidade de sobreposição regional do coeficiente Dice.
\end{itemize}